% !TEX root = ../main.tex
\chapter{Topological Manifolds}\label{ch:c02}
\begin{paracol}{2}

\begin{sources}
\begin{tabularx}{\linewidth}{@{}l l L@{}}
\toprule
\textbf{Source} & \textbf{Length} & \textbf{Content}\\
\midrule
\seg{2}{1}{Lecture 2: Topological Manifolds} & 1:23:01 & Why restrict to charted spaces; the definition of a $d$-dimensional topological manifold;
examples (torus, a wire, a non-example); charts, atlas, chart maps and coordinates; chart independence of physics; two charts of $\R^2$ minus a point; chart transition
maps and gluing; the manifold philosophy (properties defined through charts, well-definedness, why differentiability needs a smaller atlas).\\
\seg{2}{2}{Tutorial 2: Topological Manifolds} & 0:29:37 & Solutions of the sheet by a tutor (not Schuller): multiple-choice questions on manifolds, an atlas for the
M\"obius strip, the manifold glued from four charts, the curve $m^4+n^4=1$ with four charts and a transition map.\\
\bottomrule
\end{tabularx}

\smallskip
\textbf{Transcripts} \file{materials/transcripts/C02-V01.txt} and \file{C02-V02.txt} (YouTube captions).
\textbf{Sheet 2:} see \url{https://tales.mbivert.com/on-heraeus-winter-school-tutorials/}.
\end{sources}

\begin{reading}[Where this chapter sits in the other notes]
\emph{Geometric Anatomy of Theoretical Physics} notes: ``Topological manifolds and manifold bundles'' (chapter 6: charts, atlases, transition maps, the
maximal atlas and the construction of examples) and, for the invariants mentioned at the beginning (Hausdorff, connected, compact, \dots),
``Topological spaces: some heavily used invariants'' (chapter 5); ``Differentiable structures: definition and classification'' (chapter 7) continues
with the step announced at the end of this lecture (a restricted atlas).
\end{reading}

% ==================================================================
\section{Why manifolds}\label{sec:c02-why}

Last time: topological spaces \ts{2}{1}{0:00}. There are so many of them that mathematicians cannot even classify them (``that may be a statement that is a function of time'')
\ts{2}{1}{0:28}\ts{2}{1}{0:50}. There are many topological notions: Hausdorff spaces, connected spaces, paracompact spaces, compact spaces, and so on. One could take
all the notions ever invented that are properties of a topological space defined in terms of the topology, and make a big chart: this space has the notion, that one has
not; there are even notions that are not Boolean-valued but, say, group-valued. Then one asks whether this set of notions decides when two topological spaces are
\emph{homeomorphic}, i.e.\ when there is a map invertible in both directions, continuous both ways, mapping one topology onto the other. The answer: there is no known set of
topological notions with which one could decide this, since for any set of notions one finds counterexamples; so topological spaces are not classified
\ts{2}{1}{0:50}\ts{2}{1}{1:10}\ts{2}{1}{1:31}\ts{2}{1}{1:51}\ts{2}{1}{2:12}.

That is a mathematical problem. The \emph{physical} problem is that many topological spaces resemble what we think of as spacetime so little that they are not deemed
useful as models for it \ts{2}{1}{2:12}\ts{2}{1}{2:33}. For classical spacetime physics (``again a statement that may be a function of time'') one may focus on topological
spaces (recall: a pair of a set and a topology) that can be \emph{charted}, analogously to how the surface of the Earth is charted: one buys many charts that cover the surface in an
atlas in a shop, but one does not need to, since they can be constructed. These special topological spaces are called \emph{topological manifolds}
\ts{2}{1}{2:33}\ts{2}{1}{2:54}\ts{2}{1}{3:14}\ts{2}{1}{3:35}\ts{2}{1}{3:56}.

Schuller's philosophy: a bit of imprecise talk to give an idea, and then the precise definition afterwards; in the future one refers only to the precise one \ts{2}{1}{4:18}.

\section{Definition of a topological manifold}\label{sec:c02-def}

\begin{definition}[New concepts: topological manifold]\label{def:c02-manifold}
A topological space $(M,\mathcal O)$ is called a \textbf{$d$-dimensional topological manifold} if for every point $p\in M$ there exist an open set
$U\in\mathcal O$ with $p\in U$ and a map $x\colon U\to x(U)\subseteq\R^d$ (the same $d$) such that
\begin{enumerate}
  \item $x$ is invertible, i.e.\ there is an inverse map $x^{-1}\colon x(U)\to U$;
  \item $x$ is continuous;
  \item $x^{-1}$ is continuous
\end{enumerate}
\ts{2}{1}{4:18}\ts{2}{1}{4:38}\ts{2}{1}{4:58}\ts{2}{1}{5:21}\ts{2}{1}{5:42}\ts{2}{1}{6:03}\ts{2}{1}{6:24}\ts{2}{1}{6:45}\ts{2}{1}{7:07}.
\end{definition}
Remarks, as given in class:
\begin{itemize}
  \item A topological manifold is a topological space with an extra condition, the \emph{charting} condition \ts{2}{1}{4:58}.
  \item $x$ need not be surjective onto $\R^d$: it goes to its image $x(U)$, which lies in $\R^d$ \ts{2}{1}{6:03}\ts{2}{1}{6:24}.
  \item The continuity statements refer to topologies: on $U$ the subset topology inherited from $\mathcal O$ (and $U$ is open, so this is not a problem), on $\R^d$
        the standard topology, ``because nothing else is said'' \ts{2}{1}{6:45}\ts{2}{1}{7:07}\ts{2}{1}{7:26}.
  \item One might think continuity of $x$ plus invertibility already gives continuity of the inverse; the example of the last lecture ($\{1,2\}$ with two topologies) shows
        that it does not, but here it is required \ts{2}{1}{7:26}\ts{2}{1}{7:47}. It is a very simple definition.
  \item The topology $\mathcal O$ is whatever it is: it need not be the standard one (and it can't be, since if $M$ is not $\R^d$ one does not even know how to define the standard
        topology). For every point there exists an open set containing it, together with an entire map, ``a lot of data whose existence I require'', which takes the points of the
        set bijectively to a subset of $\R^d$, continuous in both directions \ts{2}{1}{7:47}\ts{2}{1}{8:10}\ts{2}{1}{8:30}\ts{2}{1}{8:51}.
  \item If $x$ is invertible with continuous inverse, then the pre-image of $U$ under $x^{-1}$ is $x(U)$; as $U$ is open, $x(U)$ is open in $\R^d$
        \ts{2}{1}{9:12}\ts{2}{1}{9:32}.
\end{itemize}

\section{Three pictures}\label{sec:c02-examples}

Before giving the elements of the definition names, Schuller gives three examples that he draws, to be calculated later, and picks up the names from navigating maps of the Earth
as the captain of a ship \ts{2}{1}{9:32}\ts{2}{1}{9:53}.

\begin{example}[The torus: a two-dimensional manifold]\label{ex:c02-torus}
Take the surface of a doughnut (of a car tyre) as the set $M$ \ts{2}{1}{9:53}\ts{2}{1}{10:13}\ts{2}{1}{10:34}. It carries a topology: it may inherit the subset topology from the
standard topology of the surrounding $\R^3$ in which one embeds it (``this is not the standard topology, but the subset topology that comes from it''), or one may prescribe one oneself
\ts{2}{1}{10:34}\ts{2}{1}{10:54}. The claim is that $M$ is a two-dimensional topological manifold (not discussed technically, but the definition is run through)
\ts{2}{1}{10:54}\ts{2}{1}{11:17}. Mark a point blindly: there exists an open subset around it; the open subsets of the induced topology look like $U'\cap M$ with $U'$ an open set of $\R^3$
(a blob without boundary), intersected with the tyre: a patch without boundary \ts{2}{1}{11:17}\ts{2}{1}{11:37}\ts{2}{1}{11:58}. Then one needs a map taking every point of this patch to some part
of $\R^2$, one to one and continuous in both directions, giving an open region in $\R^2$ (it is open because of invertibility and continuity in both directions); one can calculate and prove that
such maps exist \ts{2}{1}{11:58}\ts{2}{1}{12:19}\ts{2}{1}{12:41}\ts{2}{1}{13:02}\ts{2}{1}{13:22}.
\end{example}
This is precisely what one does for the surface of the Earth \ts{2}{1}{13:22}\ts{2}{1}{13:42}. Before the terminology, some physics \ts{2}{1}{13:42}.

\paragraph{The foot of the Eiffel Tower.} Consider the torus (``the globe is not a sphere, it could be a torus'') as the real world, and the little cross to be the foot of the Eiffel Tower, a
physically existing point you can go and touch \ts{2}{1}{13:42}\ts{2}{1}{14:04}. The map $x$ takes this point of the real world to an element of $\R^2$, which is just the set of pairs of real
numbers (``one could have chosen a totally different map''): $x(p)=(x^1(p),x^2(p))$, say $(1,2)$. Does ``the Eiffel Tower is at $(1,2)$'' mean anything? It means something only because
one has chosen a specific map: one buys an atlas, goes to page 27, and this part of $\R^2$ is page 27; one finds the Eiffel Tower and Munich; that it is at $(1,2)$ is just one centimetre
from the lower edge of the chart, because the publisher decided so; had he chosen otherwise, or distorted the chart (``a rubber paper''), things would look different: physical meaning nil.
But the foot of the Eiffel Tower in the real world is real, and we represent the real world by coordinates in a chart \ts{2}{1}{14:04}\ts{2}{1}{14:25}\ts{2}{1}{14:47}\ts{2}{1}{15:08}\ts{2}{1}{15:29}\ts{2}{1}{15:49}\ts{2}{1}{16:09}.

\begin{definition}[New concepts: chart]\label{def:c02-chart}
A pair $(U,x)$, where $x$ is a map (remember: not the coordinates, but a map taking all the points of $U$ into $\R^d$) with the properties of \cref{def:c02-manifold}, is called a
\textbf{chart} of the manifold. $U$ is the \emph{chart domain}, $x(U)$ its image \ts{2}{1}{16:30}\ts{2}{1}{16:51}.
\end{definition}
In the example: $U$ could be the area of Paris in the real world, which one can go and touch; $x(U)$ is the image of the area of Paris in the chart, which one can also touch
(the paper of the atlas); the connection is what the publisher decided to choose (he also decides whether to take Paris with or without the Bois de Boulogne [?]). The definition requires
that around \emph{every} point of the manifold there is a chart containing it: if you buy an atlas and your hometown is not on it, you return it to the shop and say you feel cheated
\ts{2}{1}{16:51}\ts{2}{1}{17:11}\ts{2}{1}{17:33}\ts{2}{1}{17:54}\ts{2}{1}{18:14}.

\begin{example}[A wire: a one-dimensional manifold]\label{ex:c02-wire}
A distorted figure-eight wire (``it is not disconnected here, it is just behind''), lying in some $\R^3$ with the subset topology of the standard topology \ts{2}{1}{18:14}\ts{2}{1}{18:34}\ts{2}{1}{18:55}.
The claim is that this is a one-dimensional topological manifold: for every point there is an open domain around it, with every point mapped one to one to some part of $\R$ (it could be
disconnected or not); the charts are one-dimensional because it is a one-dimensional world \ts{2}{1}{18:55}\ts{2}{1}{19:18}\ts{2}{1}{19:39}. (The wire is more intuitive to us although it is the same
sort of object one dimension lower.)
\end{example}

\begin{example}[A non-example: a branching line]\label{ex:c02-nonex}
A set $M$ shaped like a line that bifurcates, as a subset of $\R^2$ carrying the subset topology of the standard topology: a topological space, but \emph{not} a topological manifold
\ts{2}{1}{19:39}\ts{2}{1}{20:00}\ts{2}{1}{20:21}\ts{2}{1}{20:42}. Investigation: for the points on the branches there is no problem; at the branching point there is: an open neighbourhood exists
(take an open set in $\R^2$; its intersection with $M$ is open in the subset topology), but one cannot find from such an open set a map, invertible and continuous in both directions, into $\R^1$ or $\R^2$:
``intuitively one sees it, how would you chart this? In fact you can't.'' So it is not a ``well, question mark-dimensional'' topological manifold \ts{2}{1}{20:42}\ts{2}{1}{21:04}\ts{2}{1}{21:24}\ts{2}{1}{21:44}\ts{2}{1}{22:05}\ts{2}{1}{22:26}.
\end{example}
It is very easy to come up with something that is not a topological manifold. String theorists like to think of D-branes with little strings attached: such a configuration is not a
manifold (one-dimensional here, two-dimensional there); the definition requires one $d$ that applies everywhere. Could one loosen this? It would not solve the problem of what to do at the point
where the pieces start touching \ts{2}{1}{22:26}\ts{2}{1}{22:49}\ts{2}{1}{23:10}.

\begin{coursenote}[Why a topology had to be given]
To decide whether the branching set is a manifold, a set alone is not enough: one needs a topology. Schuller uses the standard topology of the surrounding space (restricted to $M$) because it serves our intuition for the
quickly drawn arguments, pulling back along the intersections: say one takes any open set of $\R^2$ (the shaded area, every point of which has a little soft ball inside); intersected with the subset it gives what \emph{looks}
like two open intervals. For a person coming from the standard topology it seems one already knew these to be open; but the real argument is by the subset topology \ts{2}{1}{23:10}\ts{2}{1}{23:31}\ts{2}{1}{23:52}\ts{2}{1}{24:12}\ts{2}{1}{24:33}\ts{2}{1}{24:53}\ts{2}{1}{25:13}.
\end{coursenote}

\paragraph{Student question: why is the wire one-dimensional and not two?} Could one not take $d=2$? There is no continuous map from the chart domain into $\R^2$ of the required kind: one can perhaps construct a bijective map
into $\R^2$ (already difficult, one has to wind the interval; ``I am not entirely sure, I think that's somehow possible''), but certainly not a continuous one with continuous inverse; so the definition fails if $d$ is increased to
two. ``This is very intuitive and I want to give you an impression; you need to trust me that it is all true''; specific examples will be calculated, also in the tutorials today
\ts{2}{1}{25:13}\ts{2}{1}{25:37}\ts{2}{1}{25:58}\ts{2}{1}{26:18}\ts{2}{1}{26:38}.

\section{Terminology: charts, atlas, coordinates}\label{sec:c02-terminology}

Now ``a lot of terminology, but very good terminology'' \ts{2}{1}{26:38}\ts{2}{1}{27:00}. The pair $(U,x)$ from the definition is a \emph{chart} of the manifold $(M,\mathcal O)$ (\cref{def:c02-chart}); the definition
says that around every point there is a chart. The publisher who sells an atlas must make sure of this \ts{2}{1}{27:00}\ts{2}{1}{27:20}.

\begin{definition}[New concepts: atlas]\label{def:c02-atlas}
A collection $\mathcal A=\{(U_\alpha,x_\alpha)\}_{\alpha\in A}$ of charts, with $\alpha$ a label from an arbitrary (possibly huge) index set $A$, is called an \textbf{atlas} of the topological manifold if
\[
  \bigcup_{\alpha\in A}U_\alpha=M,
\]
i.e.\ the chart domains, which may have a huge overlap, cover the entire ``real world'' \ts{2}{1}{27:20}\ts{2}{1}{27:42}\ts{2}{1}{28:03}\ts{2}{1}{28:24}\ts{2}{1}{28:44}.
\end{definition}
Unfortunately spacetime is four-dimensional: one cannot print an atlas of it on sheets of paper, one would need four-dimensional sheets, which mathematically one can do, practically one cannot; ``it is lucky that the surface of the Earth is
two-dimensional'' \ts{2}{1}{28:44}\ts{2}{1}{29:04}\ts{2}{1}{29:25}.

The map $x$ is called the \textbf{chart map} (one can have the same domain in the real world but choose different chart maps, represent it differently in the atlas) \ts{2}{1}{29:25}\ts{2}{1}{29:46}.

\begin{definition}[New concepts: coordinate maps and coordinates]\label{def:c02-coordinates}
A map $x\colon U\to\R^d$ gives to every point $p\in U$ a $d$-tuple $x(p)=(x^1(p),\dots,x^d(p))$ (recall: $\R^d$ is the set of $d$-tuples, $d$ numbers in an order), and
this information is the same as giving $d$ many maps $x^i\colon U\to\R$, $i=1,\dots,d$, where $x^1$ is a totally different map from $x^2$. These are called the \textbf{coordinate maps}
(the component maps of the chart map). For $p\in U$, the real number $x^i(p)$ is the \textbf{$i$-th coordinate of $p$ with respect to the chosen chart}
\ts{2}{1}{29:46}\ts{2}{1}{30:07}\ts{2}{1}{30:29}\ts{2}{1}{30:49}\ts{2}{1}{31:10}\ts{2}{1}{31:31}\ts{2}{1}{31:53}\ts{2}{1}{32:13}\ts{2}{1}{32:34}\ts{2}{1}{32:57}\ts{2}{1}{33:19}.
\end{definition}
With another chart this will be another number: $x^2(p)$ is what happens to be the number printed on the publisher's chart in the direction in which he puts the foot of the Eiffel Tower \ts{2}{1}{33:19}\ts{2}{1}{33:40}.

\begin{coursenote}[Chart-independent physics]
Man-made charts are a figment of our imagination: charts have nothing to do with physics; they may be very useful, but nobody would say, because a publisher printed an atlas upside down, that it is less comfortable to live in that
part of the world. Most of undergraduate physics is formulated in atlases and one very much confuses what is imagined and what is real; Schuller will therefore emphasise very often what is the real world and what is an arbitrary image of it:
``it took Einstein three years to clarify this, it takes us three days'' \ts{2}{1}{33:40}\ts{2}{1}{34:00}\ts{2}{1}{34:21}\ts{2}{1}{34:41}.\par
It is totally obvious from Lagrangian mechanics: a physical system moves on some configuration space and one can choose generalised coordinates, other generalised coordinates; choosing coordinates is choosing charts. Does the bird that flies by care about
my idea about charts? Let's take a stone, without free will: what would physics be if the stone were impressed by one's thinking about a chart? ``It would be physics with telekinesis: I could start moving things by imagining different charts.'' Frames of reference,
generalised coordinates \ts{2}{1}{34:41}\ts{2}{1}{35:02}\ts{2}{1}{35:23}\ts{2}{1}{35:43}. If an experiment shows that by thinking one can move things, we do chart-dependent physics; since that is not the case, all we do until further notice is
\emph{chart-independent} physics. The foot of the Eiffel Tower is a real point; changing coordinates does not move it: ``we do physics without telekinesis, that is the programmatic aim'' \ts{2}{1}{35:43}\ts{2}{1}{36:05}\ts{2}{1}{36:25}\ts{2}{1}{36:46}.
\end{coursenote}

\section{An example: \texorpdfstring{$\R^2$}{R2} without a point, two charts}\label{sec:c02-r2}

Take as the real world the points of $\R^2$, all the pairs. The chart domain $U$ is $\R^2$ without the element $(0,0)$ \ts{2}{1}{36:46}\ts{2}{1}{37:06}\ts{2}{1}{37:26}.

\paragraph{An aside: a set versus a topological space.} What is the difference between the set $\{1,2\}$ and the set $\{2,1\}$? None, a set has no order. $\R^2$ as a set is like a \emph{powder}: ``I can rip it apart, I can grind it, I can make a powder of this, fill
the powder in a glass, shake the glass, and pour the little pile of powder in front of you and say: there you got your $\R^2$.'' Shaking changes all relations between the points one may imagine. Equipped with the standard topology it is more like a rubber sheet: one can pull it, but not
rip it apart, because that would be discontinuous. The $\R^2$ one usually thinks of is the topological space with the standard topology, but as a pure set it is a powder; a warning about what a set, a topological space, a topological manifold is
\ts{2}{1}{37:26}\ts{2}{1}{37:48}\ts{2}{1}{38:08}\ts{2}{1}{38:29}\ts{2}{1}{38:50}\ts{2}{1}{39:10}.

\begin{example}[Two charts of $\R^2\setminus\{(0,0)\}$]\label{ex:c02-two-charts}
\textbf{Chart $x$.} Let $x\colon U\to\R^2$, $(m,n)\mapsto(-m,-n)$ (an arbitrary chart invented ``because I feel like it''; it is an element of $\R^2$; the origin was taken out ``for fun'')
\ts{2}{1}{39:10}\ts{2}{1}{39:31}\ts{2}{1}{39:51}. Picture: the real blackboard, without a point, is mapped to a chart, a picture of the blackboard as in one's book: every point goes to its opposite, chalk
dust is mapped into the chart point by point; ``you might say it inverts the real world: what do you mean, it is a chart map, it is as good as any other'' \ts{2}{1}{39:51}\ts{2}{1}{40:11}\ts{2}{1}{40:32}\ts{2}{1}{40:53}\ts{2}{1}{41:13}.\\
\textbf{Chart $y$.} On the same set (more precisely on the slit plane) take
\[
  y\colon(m,n)\longmapsto\Bigl(\sqrt{m^2+n^2},\ \arctan\tfrac nm\Bigr)
\]
(the angle taken in $(0,2\pi)$; the formula is schematic [?])
``if you think that is crazy, then you say polar coordinates are crazy'': the first coordinate is the distance from the origin, the second an angle with, say, the $m$-axis. The chart domain is the
\emph{slit plane}: one takes out not only the point but the entire half axis $\{(a,0)\mid a\in\R_0^+\}$ (including the origin; the negative half axis is still in), ``for certain reasons, do not worry'';
the image is a half-infinite strip $(0,\infty)\times(0,2\pi)$ [?]. A circle of the real world appears in this chart as a segment of a straight line (the point where the slit cuts it is missing, being not in the domain)
\ts{2}{1}{42:56}\ts{2}{1}{43:17}\ts{2}{1}{43:38}\ts{2}{1}{43:58}\ts{2}{1}{44:19}\ts{2}{1}{44:39}\ts{2}{1}{45:01}\ts{2}{1}{45:24}\ts{2}{1}{45:45}\ts{2}{1}{46:08}.
\end{example}
This is what is called generalised coordinates in classical mechanics ($r$ and $\varphi$): here they are $y^1(p)$ and $y^2(p)$, ``this is just what it is''. The same real-world situation looks very different in the two charts (in the first the
circle remains a circle, you just do not see that it is inverted); any chart one can come up with is covered by the definition, and one needs to know how the charts are related to the real world
\ts{2}{1}{46:08}\ts{2}{1}{46:29}\ts{2}{1}{46:49}\ts{2}{1}{47:09}.

\paragraph{Student question: is the real world not already $\R^2$?} A student objects that taking the points of $\R^2$ as the ``real world'' already assumes numbers. ``You are absolutely right, it is a notorious problem of emphasising'': the real world does not
consist of numbers. What would be needed is what was done in the beginning: point with a finger at an element of the blackboard and assign to it the number $(-1,-2)$; in real life one puts little flags, ``here on the Moon is our flag and we call this point $(1,2)$'';
that is inventing a chart, and it is what gold diggers do, putting down their claims and drawing lines. In physics we always ignore the real world and talk only about charts and chart representatives of it, but at least mathematically one should put the real world in to know what one is talking about
\ts{2}{1}{41:13}\ts{2}{1}{41:34}\ts{2}{1}{41:55}\ts{2}{1}{42:15}\ts{2}{1}{42:35}\ts{2}{1}{42:56}.

\section{Chart transition maps}\label{sec:c02-transition}

``These are derived objects, objects derived from charts we already have, so actually not worth mentioning, but they are heavily used by physicists, hence we introduce them'' \ts{2}{1}{47:09}\ts{2}{1}{47:36}.
Imagine two charts $(U,x)$ and $(V,y)$ with different domains and different prescriptions, whose domains overlap: $U\cap V\neq\emptyset$. A point of the overlap (the foot of the Eiffel Tower) can be mapped either by $x$ (as part of $U$)
or by $y$ (as part of $V$): ``you turn the pages of your atlas and the Eiffel Tower appears more than once, on this chart and on the other; the charts always have a little overlap in order for you to navigate'' \ts{2}{1}{47:36}\ts{2}{1}{48:01}\ts{2}{1}{48:22}\ts{2}{1}{48:44}\ts{2}{1}{49:05}\ts{2}{1}{49:26}.
Both maps can be restricted to $U\cap V$, giving $x(U\cap V)$ and $y(U\cap V)$ \ts{2}{1}{49:47}\ts{2}{1}{50:09}\ts{2}{1}{50:29}.

On the chart of page 7 the foot of the Eiffel Tower is here, on page 28 there: can one translate coordinates directly, without the real world? This was Gauss's concern (putting flags on the Earth to make a map)
\ts{2}{1}{50:29}\ts{2}{1}{50:50}\ts{2}{1}{51:10}. It is very simple: $x$ is invertible, so take the point in the first chart, feed it into $x^{-1}$, which leads to the real foot of the Eiffel Tower, then apply $y$ \ts{2}{1}{51:10}\ts{2}{1}{51:30}\ts{2}{1}{51:53}.

\begin{definition}[New concepts: chart transition map]\label{def:c02-transition}
For charts $(U,x)$ and $(V,y)$ the \textbf{chart transition map} (change of chart) is
\[
  y\circ x^{-1}\colon x(U\cap V)\longrightarrow y(U\cap V),
\]
a map from a subset of $\R^d$ to another subset of $\R^d$ \ts{2}{1}{52:13}\ts{2}{1}{52:34}\ts{2}{1}{52:55}\ts{2}{1}{53:16}.
\end{definition}
Since the chart maps are invertible and continuous in both directions, the transition map is a map between \emph{open} parts of $\R^d$, and it is \emph{continuous}, because it is the composition of continuous maps (last lecture)
\ts{2}{1}{53:36}\ts{2}{1}{53:58}\ts{2}{1}{54:18}.

When physicists talk about a change of coordinates or of chart, they go from one favourite imagination to another, e.g.\ from Cartesian to polar coordinates, back and forth, and never touch as a mathematical object the real world: only
coordinates of the foot of the Eiffel Tower. This is fair and correct in mechanics, but it can lead to the philosophical problem of imagining some reality to the coordinates; even after understanding that one can choose others, one might hold that the Cartesian ones are the ``real'' ones
(hence not called generalised) and the others somehow less real: this means nothing. The theory is usually presented within the yellow box of coordinates; one advantage of the manifold picture is that one can calculate very nicely there, but one must not forget the real world one wants to talk about:
the foot of the Eiffel Tower is the point $p$, and it does not care whether one looks at this chart map or that one. ``This took Einstein three years to fully understand; he understood it half, or 3/4, or 99\%, and that is wrong: you need to understand it 100\%. And this manifold picture gives you
this 100\%'' \ts{2}{1}{54:18}\ts{2}{1}{54:40}\ts{2}{1}{55:00}\ts{2}{1}{55:21}\ts{2}{1}{55:42}\ts{2}{1}{56:02}.

\paragraph{Transition maps glue the atlas together.} The chart transition maps contain fantastic information \ts{2}{1}{56:02}\ts{2}{1}{56:23}. In an atlas with many pages, one continent looks like this on one page and like that on another; one realises that this one is actually the same as that one;
``Madrid is to be found again here'', so there needs to be a transition map from here to there: informally, it contains the \emph{instructions of how to glue together the charts} of an atlas \ts{2}{1}{56:23}\ts{2}{1}{56:43}\ts{2}{1}{57:04}.

\begin{example}[The captain who glues his atlas]\label{ex:c02-glue}
Imagine the captain of a modern ship, deep down with no windows, on a long cruise with nothing to do, with a single-sided atlas of the world, desperate. He pulls out the charts and some glue and glues the charts together in their overlap
regions, which he can clearly identify. What does he get? A paper model of the sphere (a little wiggly if the charts are small enough). All he used is the information that this part looks like that part: green continent glued with green continent; he identified a point of one chart with the
corresponding point of the other and glued. One may have to bend the paper; then one gets a sphere \ts{2}{1}{57:04}\ts{2}{1}{57:27}\ts{2}{1}{57:47}\ts{2}{1}{58:08}\ts{2}{1}{58:28}\ts{2}{1}{58:49}.
\end{example}
In today's tutorial one will have charts of a world and glue them together; ``you will be surprised what you get'' \ts{2}{1}{58:49}. All the information about the real world, globally and topologically, is in the chart transition maps, hence they are beloved by physicists: a good
data structure, $\R^d\to\R^d$, very simple to implement, but the meaning is only what you glue together; ``nobody would say that, having an atlas, you understood that the surface of the Earth is a sphere, unless you start gluing'' \ts{2}{1}{58:49}\ts{2}{1}{59:09}\ts{2}{1}{59:29}.

\section{Manifold philosophy}\label{sec:c02-philosophy}

Now a very important point: the whole philosophy of writing down manifolds, of which the full extent will only be understood when differentiability of curves and functions is discussed. So far only continuity (no jumps); differentiability of a curve means that there is
everywhere a tangent, roughly speaking \ts{2}{1}{59:29}\ts{2}{1}{59:50}\ts{2}{1}{1:00:11}.

The basic idea: with an atlas one can represent the real world in charts. For instance the course of a ship in the real world (the bubbles it leaves in the sea) can be drawn in a chart; this is what navigators do. Would the navigator, or his supervisor at home, be surprised if the course of the ship
jumped to another point and continued from there? Yes, very much. What would the supervisor's argument be? That the chart maps from the real world to the chart are bijective and continuous, that in the chart one chose the standard topology, and so on; so a discontinuity of the curve in the chart, which is a curve from the
real numbers (the curve parameter) into $\R^2$, signals a discontinuity in the real world. This is what undergraduates are taught: a point particle is a curve in $\R^3$ or $\R^2$; but this is just in the chart. The curve in $\R^3$ is not the real curve in the real world, it is the image of the curve in a chart one chose
\ts{2}{1}{1:00:11}\ts{2}{1}{1:00:31}\ts{2}{1}{1:00:52}\ts{2}{1}{1:01:12}\ts{2}{1}{1:01:32}\ts{2}{1}{1:01:53}.

Because of the \emph{postulated} continuity of chart maps one can judge the continuity of the curve in the real world from continuity in the chart. One could instead define continuity directly in the real world (one can, since the real world is a topological space), but one may \emph{push} the real-world objects down into the charts,
look at their chart representatives using the chart maps, and study them in the charts, which is easy since these are curves $\R\to\R^2$ or functions $\R^2\to\R$ (like a temperature distribution drawn in a chart): objects from undergraduate analysis. Then one gets the idea of defining properties of real-world objects by
properties of their chart representatives. ``The answer is yes, you can, but you have to be careful'' \ts{2}{1}{1:01:53}\ts{2}{1}{1:02:14}\ts{2}{1}{1:02:34}\ts{2}{1}{1:02:54}.

\begin{remark}[Manifold philosophy]\label{rem:c02-philosophy}
Often it is desirable (and indeed in some cases the only way) to define properties of real-world objects by imposing suitable conditions not on the object itself but on a \emph{chart representative} or image of it
\ts{2}{1}{1:03:15}\ts{2}{1}{1:03:38}\ts{2}{1}{1:03:59}\ts{2}{1}{1:04:19}\ts{2}{1}{1:04:40}\ts{2}{1}{1:05:00}.
\end{remark}
Example: continuity of a real-world curve $\gamma\colon\R\to M$ (the course of the ship; ``a curve from some parameter set into $M$''). Take a chart $(U,x)$ and look at the portion of the curve that lies in $U$ (to make the argument simpler; it can be generalised).
Map it into the chart: the \emph{chart representative} is $x\circ\gamma$, first $\gamma$ then $x$; it is a curve from a real parameter to a part of $\R^2$: what one is told in undergraduate physics is the trajectory of a particle. In fact it is only the chart representative of
the curve in the real world, which is $\gamma$. By ``tacit assumption'' one says: if the chart representative is (undergraduate-)continuous, which is known from analysis (continuity of $\R\to\R^d$), then the particle path in the real world is continuous.
One defines a property like continuity of a real-world object by judging suitable conditions on the chart representative \ts{2}{1}{1:05:00}\ts{2}{1}{1:05:21}\ts{2}{1}{1:05:41}\ts{2}{1}{1:06:02}\ts{2}{1}{1:06:22}\ts{2}{1}{1:06:42}\ts{2}{1}{1:07:03}\ts{2}{1}{1:07:24}.

\paragraph{Advantages and disadvantages.} Advantage: undergraduate analysis can be used in charts; mathematics in charts is mathematics on the real numbers, not so complicated \ts{2}{1}{1:07:24}\ts{2}{1}{1:07:45}.
\emph{Philosophical} disadvantage: one might start confusing what is a chart image and what is the real-world object (``assuming you are Einstein, and spend three years thinking about it, you no longer fall into that trap''). A further disadvantage:
the notion defined by referring to the chart representation may be \emph{ill defined}. A definition says ``this is defined like so''; but here a property is assigned to $\gamma$ by judging the object $x\circ\gamma$, where $x$ is a figment of our imagination: one assigns to the path of the real bird a property that
also depends on our thinking. That is very dangerous unless one ensures that for any other thinking the result does not change. ``An arbitrarily chosen chart is employed; we need to make sure that the defined property does not change if we afford another fantasy, and the fantasy is another chart''
\ts{2}{1}{1:07:45}\ts{2}{1}{1:08:06}\ts{2}{1}{1:08:27}\ts{2}{1}{1:08:48}\ts{2}{1}{1:09:08}\ts{2}{1}{1:09:29}\ts{2}{1}{1:09:51}\ts{2}{1}{1:10:14}\ts{2}{1}{1:10:37}.

Formally: a curve (the same for a function on a domain, or anything) is composed with a chart map (``fantasy, arbitrary choice, superstition'') into a composite object partly of the real world; if some property holds for $x\circ\gamma$ it must hold for $y\circ\gamma$ for any other chart $(V,y)$.
If one can ensure that if the property holds in some chart it holds in any other chart, it is independent of the choice of fantasy. One knows this from equivalence relations: an operation defined on equivalence classes via representatives must be independent of the representative chosen
\ts{2}{1}{1:10:37}\ts{2}{1}{1:10:58}\ts{2}{1}{1:11:19}\ts{2}{1}{1:11:39}\ts{2}{1}{1:11:59}\ts{2}{1}{1:12:20}.

\begin{proposition}[Chart independence of continuity of a curve]\label{prop:c02-curve-indep}
Let $\gamma$ be a curve in $M$ and $(U,x)$, $(V,y)$ charts. If $x\circ\gamma$ is (undergraduate-)continuous, so is $y\circ\gamma$ (on the common domain).
\end{proposition}
\begin{proof}
$y\circ\gamma=(y\circ x^{-1})\circ(x\circ\gamma)$: the transition map $y\circ x^{-1}$ is continuous (a composition of continuous maps, last lecture), and $x\circ\gamma$ is continuous by assumption; the composition of continuous maps is continuous
\ts{2}{1}{1:12:20}\ts{2}{1}{1:12:40}\ts{2}{1}{1:13:00}\ts{2}{1}{1:13:22}\ts{2}{1}{1:13:42}\ts{2}{1}{1:14:03}. (The composition is associative; one has just inserted the identity $x^{-1}\circ x$, as a student remarks.)
\end{proof}
So the continuity of the chart transition map ``saved the day'': it makes this chart-dependent definition independent of the choice of chart \ts{2}{1}{1:14:03}\ts{2}{1}{1:14:24}.

\paragraph{Why differentiability needs more.} Now imagine that one defines $\gamma$ to be \emph{differentiable} if $x\circ\gamma\colon\R\to\R^d$ is a differentiable map (finding the derivative, which at a point one would call the velocity). One cannot decide this at the level of
the manifold, since from $\R$ into a set with a topology there is not enough structure (one cannot do sums, subtractions and divisions needed for a differential quotient); one can only do it in the chart, where $\R^d$ has addition and multiplication by scalars. But is the property independent of the chart?
No: the transition map $y\circ x^{-1}$ was only guaranteed to be continuous; composing a differentiable function with a merely continuous function may give a merely continuous one: ``the chart transition map can introduce into a nice smooth curve a little edge'', with no unique tangent
\ts{2}{1}{1:14:24}\ts{2}{1}{1:14:45}\ts{2}{1}{1:15:06}\ts{2}{1}{1:15:26}\ts{2}{1}{1:15:47}\ts{2}{1}{1:16:08}\ts{2}{1}{1:16:29}\ts{2}{1}{1:16:49}\ts{2}{1}{1:17:09}.
\emph{So on a topological manifold with merely continuous transition maps one cannot define differentiability through the charts}, and the charts are the only possibility.

How to lift the notion in the chart to the manifold, for any chart? The charts come from a big atlas. What if one takes a huge atlas with lots of topological charts with continuous transition functions and rips off pages 7, 113 and 1080 because those were the ones whose transition
functions were continuous but not differentiable? Taking out all the troublemakers gives a restricted atlas; if one agrees to take charts only from it, all the transition functions are smooth, compositions of differentiable maps are differentiable, and the problem is gone
\ts{2}{1}{1:17:09}\ts{2}{1}{1:17:31}\ts{2}{1}{1:17:52}\ts{2}{1}{1:18:12}\ts{2}{1}{1:18:33}.

\begin{definition}[New concepts: maximal atlas, and the idea of a differentiable atlas]\label{def:c02-maximal}
A \textbf{maximal (topological) atlas} is one that contains every chart it could possibly contain, ``lots of redundancy'': the maximum atlas of a topological space \ts{2}{1}{1:18:12}\ts{2}{1}{1:18:33}.
\end{definition}
Ripping out charts from a maximal atlas is giving \emph{information} (``I could encode a message in the numbers of the pages I rip out: 7, 13, \dots, the numbers in the alphabet, `I love you'''): reducing a maximal atlas by ripping out the charts that make trouble for the notions we want to define provides
additional information, and this is called the prescription of a \textbf{differentiable atlas}. That is the topic of the next lecture but one, where differentiability will be made to work, a major step in understanding spacetime
\ts{2}{1}{1:18:33}\ts{2}{1}{1:18:53}\ts{2}{1}{1:19:14}\ts{2}{1}{1:19:36}.

\paragraph{Student questions at the end.} A student (``one of those happy physicists who don't care about mathematical rigour, but since you forced me to think'') asks how the chart notion of continuity compares with the manifold notion defined in lecture 1: that $x\circ\gamma$ is undergraduate-continuous gives only
a notion of ``lifted'' continuity; has it been proven to be the continuity defined before? Not by the argument above, but it can be proven simply, in a few lines: if $x\circ\gamma$ and $x^{-1}$ are continuous then $\gamma=x^{-1}\circ(x\circ\gamma)$ is continuous in the original sense
(composition of continuous maps). Schuller de-emphasises this point: for continuity it is ``an accident of topology'' that one can decide at the manifold level at all; for differentiability there is no manifold-level notion to compare with, and one is thrown back onto the charts
\ts{2}{1}{1:19:36}\ts{2}{1}{1:19:56}\ts{2}{1}{1:20:17}\ts{2}{1}{1:20:37}\ts{2}{1}{1:20:57}\ts{2}{1}{1:21:18}\ts{2}{1}{1:21:38}\ts{2}{1}{1:21:58}.
A second student points out that for this one needs the assumption that the chart map is invertible and its inverse is continuous too. Schuller: ``I am not sure; it could be that you are right that you need the continuity of the inverse as well for the implication; I think not but I am not sure,
we do this in the break'' \ts{2}{1}{1:21:58}\ts{2}{1}{1:22:20}\ts{2}{1}{1:22:40}. Then the tutorials at 2:00.
\begin{coursenote}[The inverse is needed]
In the direction $\gamma=x^{-1}\circ(x\circ\gamma)$ the continuity of $x^{-1}$ is indeed used, so the student is right; in the other direction ($\gamma$ continuous $\Rightarrow$ $x\circ\gamma$ continuous) only the continuity of $x$ is used. (Not said in class.)
\end{coursenote}

\begin{history}[Atlas and Gauss]
Schuller names Gauss's concern with mapping the Earth's surface by putting flags on it. \emph{Checked.} Gauss's \emph{Disquisitiones generales circa superficies curvas} (1827) studies curved surfaces intrinsically, using coordinates on them (and the theorem on the invariance of curvature under bending); the abstract notion of a manifold came with Riemann's 1854 habilitation lecture, and Weyl's
\emph{Die Idee der Riemannschen Fl\"ache} (1913) gave the definition by charts and transition maps. (Not said in class.)
\end{history}

% ==================================================================
\section{Tutorial 2}\label{sec:t02-1}

Tutorial 2, 0:29:37, \seg{2}{2}{video}, by a tutor, not Schuller. The sheet (Sheet 2) is paraphrased, not copied; see \url{https://tales.mbivert.com/on-heraeus-winter-school-tutorials/} (and \url{https://tales.mbivert.com/heraeus/Sheet_02.pdf}).
The first exercise is again multiple choice with three groups of questions: tick the correct statements \ts{2}{2}{0:00}\ts{2}{2}{0:20}.

\subsection*{Exercise 1: multiple choice}

\paragraph{(A) Topological manifolds\dots}
\begin{itemize}
  \item \dots have an integer dimension: \textbf{true} \ts{2}{2}{0:20}\ts{2}{2}{0:40}.
  \item \dots are equipped with a topology and exactly one chart: \textbf{false}; they can have more than one chart (an atlas) \ts{2}{2}{0:40}.
  \item \dots are a special case of topological spaces: \textbf{true}; every manifold is a topological space, but not every topological space is a manifold \ts{2}{2}{1:03}\ts{2}{2}{1:24}.
  \item \dots topological spaces are the minimal structure to define continuity of maps (as a statement about manifolds): \textbf{false}; continuity can already be checked on topological spaces \ts{2}{2}{1:24}\ts{2}{2}{1:45}.
  \item \dots are homeomorphic to $\R^d$: \textbf{false}. Definition: $M$ and $N$ are \emph{homeomorphic} if there is a bijection $x\colon M\to N$ such that $x$ and its inverse are both continuous \ts{2}{2}{1:45}\ts{2}{2}{2:12}\ts{2}{2}{2:36}\ts{2}{2}{2:58}.
        A manifold is only \emph{locally} homeomorphic to $\R^d$: each chart domain (an open subset) is homeomorphic to a part of $\R^d$, but the manifold is usually not homeomorphic to all of $\R^d$ \ts{2}{2}{2:58}\ts{2}{2}{3:20}\ts{2}{2}{3:41}.
\end{itemize}

\paragraph{(B) Which statements are correct?}
\begin{itemize}
  \item Every topological manifold has a maximal atlas: \textbf{true}; it comes with some atlas, and a theorem says that from any atlas one can uniquely build a maximal one \ts{2}{2}{4:03}.
  \item A maximal atlas may contain infinitely many charts: \textbf{true} \ts{2}{2}{4:23}.
  \item Continuity of a curve on a manifold can be checked with respect to an atlas: \textbf{true} (seen in the lecture) \ts{2}{2}{4:23}\ts{2}{2}{4:46}.
  \item An atlas can never contain only one chart: \textbf{false}, nonsense; an example will come later \ts{2}{2}{4:46}\ts{2}{2}{5:06}.
  \item If a function is continuous in one chart, it is continuous in every chart of a maximal atlas: \textbf{true} (the question is a little ambiguous): the function is continuous in another chart with overlap by the composition of continuous maps; one continues with further overlapping charts, so it
        is continuous in all charts of a maximal atlas, restricting to charts that overlap \ts{2}{2}{5:06}\ts{2}{2}{5:32}\ts{2}{2}{5:52}\ts{2}{2}{6:13}.
\end{itemize}

\paragraph{(C) Further statements.}
\begin{itemize}
  \item $\R$ with the standard topology and the atlas whose only chart is the identity map is a one-dimensional topological manifold: \textbf{true}, an example of an atlas with exactly one chart \ts{2}{2}{6:13}\ts{2}{2}{6:34}\ts{2}{2}{6:55}.
  \item A topological manifold can have a finite underlying set: \textbf{false}; it is locally homeomorphic to parts of $\R^d$, which are uncountable, and a finite set cannot look like something uncountable \ts{2}{2}{6:55}\ts{2}{2}{7:16}\ts{2}{2}{7:37}.
  \item A topological manifold can never have a subset topology: \textbf{false}, nonsense; many examples in the lecture had one \ts{2}{2}{7:37}\ts{2}{2}{7:58}.
  \item On the $42$-dimensional sphere one cannot build an atlas: \textbf{false}, there is no reason why one could not \ts{2}{2}{7:58}\ts{2}{2}{8:19}.
  \item A function $M\to\R$ can only be called continuous if $M$ is a manifold: \textbf{false}; continuity can already be checked on a topological space \ts{2}{2}{8:19}\ts{2}{2}{8:40}.
\end{itemize}

\subsection*{Exercise 2: an atlas for the M\"obius strip}

A M\"obius strip printed with a river: it can be decomposed on paper as a rectangle whose one edge (running from top to bottom) is identified with the opposite edge, but the other way round (from bottom to top) \ts{2}{2}{8:40}\ts{2}{2}{9:01}\ts{2}{2}{9:22}\ts{2}{2}{9:46}\ts{2}{2}{10:09}\ts{2}{2}{10:30}.
It cannot be covered with one chart; choose two charts with domains $U_1$ and $U_2$ that overlap in two places. The chart $x$ maps $U_1$ to a part of $\R^2$, where the strip with its river looks, say, like a piece of rectangle with a river across; the chart $y$ maps $U_2$ to another part of $\R^2$ where the river looks
different. The two charts form an atlas of the M\"obius strip with the river \ts{2}{2}{10:30}\ts{2}{2}{10:54}\ts{2}{2}{11:18}\ts{2}{2}{11:43}\ts{2}{2}{12:05}\ts{2}{2}{12:26}\ts{2}{2}{12:47}\ts{2}{2}{13:12}\ts{2}{2}{13:33}.

\subsection*{Exercise 3: from an atlas to a manifold}

The other way round: one is given four charts (pictures of rectangles) and has to build the manifold. One finds the overlaps of the four charts and sees them overlap at the bottom and the top, and on the left and the right: so one edge is identified with the opposite edge in the \emph{same} orientation, and likewise
bottom and top in the same direction: the result is a \textbf{torus} (the captions are garbled here [?], the picture is the torus), equipped with these charts. So a manifold has been built from an atlas, as before an atlas from a manifold \ts{2}{2}{13:33}\ts{2}{2}{13:54}\ts{2}{2}{14:27}\ts{2}{2}{14:51}\ts{2}{2}{15:16}\ts{2}{2}{15:41}\ts{2}{2}{16:01}\ts{2}{2}{16:22}.

\subsection*{Exercise 4: the curve m4+n4=1}

Let $F_1=\{(m,n)\in\R^2\mid m^4+n^4=1\}$ with the subset topology of the standard topology of $\R^2$ \ts{2}{2}{16:22}\ts{2}{2}{16:48}\ts{2}{2}{17:08}.
\begin{enumerate}
  \item The map $x\colon F_1\to[-1,1]$, $(m,n)\mapsto m$ is \emph{not injective}: $(0,1)$ and $(0,-1)$ both go to $0$ \ts{2}{2}{17:08}\ts{2}{2}{17:29}\ts{2}{2}{17:52}\ts{2}{2}{18:12}\ts{2}{2}{18:33}\ts{2}{2}{18:53}.
  \item It becomes injective on either of two maximal open subsets: $x^\uparrow\colon\{(m,n)\in F_1\mid n>0\}\to(-1,1)$, $(m,n)\mapsto m$ (strict inequality so that the domain is open and the target the open interval), and
        $x^\downarrow$ on $\{n<0\}$, again to $(-1,1)$ \ts{2}{2}{18:53}\ts{2}{2}{19:14}\ts{2}{2}{19:36}\ts{2}{2}{20:00}\ts{2}{2}{20:21}\ts{2}{2}{20:43}\ts{2}{2}{21:04}.
  \item Likewise $y^\uparrow\colon\{(m,n)\in F_1\mid m>0\}\to(-1,1)$, $(m,n)\mapsto n$ is injective, and surjective, hence bijective; its inverse is
        \[
          (y^\uparrow)^{-1}(a)=\bigl(\sqrt[4]{1-a^4},\,a\bigr),\qquad a\in(-1,1),
        \]
        \ts{2}{2}{21:04}\ts{2}{2}{21:27}\ts{2}{2}{21:48}\ts{2}{2}{22:09}\ts{2}{2}{22:30}\ts{2}{2}{22:51}\ts{2}{2}{23:11}\ts{2}{2}{23:33}\ts{2}{2}{23:55}\ts{2}{2}{24:15}, and similarly $y^\downarrow$ on $\{m<0\}$.
  \item The domains of $x^\uparrow$ and $y^\uparrow$ overlap in $\{(m,n)\in F_1\mid m>0,\ n>0\}$ \ts{2}{2}{24:36}\ts{2}{2}{24:58}. The transition map $x^\uparrow\circ(y^\uparrow)^{-1}$ is defined on the image of the overlap under $y^\uparrow$, the interval $(0,1)$, and maps
        $s\mapsto x^\uparrow\bigl(\sqrt[4]{1-s^4},s\bigr)=\sqrt[4]{1-s^4}\in(0,1)$ \ts{2}{2}{24:58}\ts{2}{2}{25:19}\ts{2}{2}{25:40}\ts{2}{2}{26:02}\ts{2}{2}{26:23}\ts{2}{2}{26:47}\ts{2}{2}{27:07}.
  \item How many charts cover $F_1$? $x^\uparrow$ and $x^\downarrow$ cover $F_1$ without the points $(\pm1,0)$; $y^\uparrow$ and $y^\downarrow$ cover it without the points $(0,\pm1)$. So four charts are needed, and their domains and maps form an atlas: $F_1$ is a one-dimensional topological manifold
        \ts{2}{2}{27:07}\ts{2}{2}{27:28}\ts{2}{2}{27:49}\ts{2}{2}{28:09}\ts{2}{2}{28:31}\ts{2}{2}{28:54}\ts{2}{2}{29:15}.
\end{enumerate}

% ==================================================================
\section*{Key concepts}
\begin{keyconcepts}
\textbf{Retrieval questions.}
\begin{enumerate}
  \item Why do we restrict from topological spaces to topological manifolds for spacetime?
  \item State the definition of a $d$-dimensional topological manifold. Where is the continuity of the inverse used?
  \item Give an example and a non-example of a manifold, and say why the latter fails at one point.
  \item What are a chart, an atlas, the chart map, the coordinate maps and the coordinates of a point?
  \item Why is changing coordinates not changing the foot of the Eiffel Tower? What would chart-dependent physics be?
  \item Define the chart transition map. Why is it continuous, and what does it say about gluing?
  \item What is the manifold philosophy, and what can go wrong with it (ill-defined properties)?
  \item Why can differentiability not be defined with just the maximal topological atlas?
\end{enumerate}
\textbf{Concepts.}
\begin{itemize}
  \item $d$-dimensional topological manifold: every point has an open neighbourhood $U$ and a homeomorphism $x\colon U\to x(U)\subseteq\R^d$ (\cref{def:c02-manifold}); chart $(U,x)$ (\cref{def:c02-chart}).
  \item Atlas: charts whose domains cover $M$ (\cref{def:c02-atlas}); coordinate maps $x^i$ and coordinates (\cref{def:c02-coordinates}).
  \item Chart transition map $y\circ x^{-1}\colon x(U\cap V)\to y(U\cap V)$, continuous (\cref{def:c02-transition}); it encodes how to glue the charts (\cref{ex:c02-glue}).
  \item Manifold philosophy: define properties via chart representatives $x\circ\gamma$ and check independence of the chart (\cref{rem:c02-philosophy,prop:c02-curve-indep}).
  \item Maximal atlas, and restricted (differentiable) atlas as information (\cref{def:c02-maximal}).
  \item Non-examples: branching line (\cref{ex:c02-nonex}); a finite set is never a manifold.
\end{itemize}
\end{keyconcepts}

\flushside
\end{paracol}
